% ---
% filename : ["electrotechnique_puissance_20260429_842_intensite_disjoncteur_resistance.tex"]
% id : "electrotechnique_puissance_20260429_842"
% title : "Intensité, disjoncteur et résistance d'un appareil"
% domain : "Électrotechnique"
% subdomain : "Puissance"
% tags : ["énergie_électrique", "intensité", "disjoncteur", "résistance", "loi_ohm", "ts"]
% difficulty : 1
% time_solve : 6
% has_octave : true
% octave_files : ["electrotechnique_puissance_20260429_842_intensite_disjoncteur_resistance.m"]
% author : "om"
% ---
\documentclass[a4paper,11pt,answers,addpoints]{exam}

% ---------------------------------------------------------
% LANGUE, ENCODAGE, FONTS
% ---------------------------------------------------------
\usepackage[utf8]{inputenc}

% ---------------------------------------------------------
% GÉOMÉTRIE ET MISE EN PAGE
% ---------------------------------------------------------
\usepackage[left=1.7cm,right=1.5cm,top=2cm,bottom=1cm]{geometry}
%\usepackage{a4wide}
\usepackage{microtype}      % Meilleure répartition du texte
\usepackage{float}          % Positionnement [H]
\usepackage{needspace}      % Saut conditionnel intelligent

% Éviter les veuves/orphelines (optionnel, à décommenter si besoin)
% \widowpenalty=10000
% \clubpenalty=10000

% Espacement vertical flexible
%\raggedbottom
%\setlength{\parskip}{0.5em plus 0.2em minus 0.1em}
%\setlength{\parindent}{0pt}


% ---------------------------------------------------------
% MATHÉMATIQUES
% ---------------------------------------------------------
\usepackage{amsmath, amssymb, bm, esvect}

% Vecteurs
\newcommand{\V}{\overrightarrow}
\def\vect#1{\overrightarrow{\strut#1}}
\providecommand{\Degrees}[0]{\ensuremath{^\circ}}

% ---------------------------------------------------------
% UNITÉS ET GRANDEURS (SIunitx)
% ---------------------------------------------------------
\usepackage{siunitx}
\sisetup{output-decimal-marker={,}}
\DeclareSIUnit \var {var}
\DeclareSIUnit \VA {VA}
\DeclareSIUnit \kVA {kVA}
\DeclareSIUnit[quantity-product={}] \degree{\text{\textdegree}}

% ---------------------------------------------------------
% GRAPHIQUES : TikZ + CircuitikZ + PGFPlots
% ---------------------------------------------------------
\usepackage{tikz}
\usetikzlibrary{
	arrows.meta,
	intersections,
	decorations.pathreplacing,
	calligraphy,
	positioning
}

\usepackage[european,straightvoltages,EFvoltages]{circuitikz}

\usepackage{tkz-fct}
\usepackage{pgfplots}
\pgfplotsset{compat=1.12}

% Symbole étoile-triangle personnalisé
\newcommand{\wye}{%
	\mathbin{\tikz[x=1ex,y=1ex]{%
			\draw[line width=.1ex] (0,0)--(45:1)--++(-45:1) (45:1)--++(0,1);
	}}%
}


\usepackage{sankey}

\pgfdeclarelayer{background}
\pgfdeclarelayer{foreground}
\pgfdeclarelayer{sankeydebug}
\pgfsetlayers{background,main,foreground,sankeydebug}


% ---------------------------------------------------------
% TABLEAUX
% ---------------------------------------------------------
\usepackage{tabularx}
\usepackage{tabu}

% ---------------------------------------------------------
% HYPERLIENS
% ---------------------------------------------------------
\usepackage[colorlinks=true,
menucolor=red,
breaklinks=true,
urlcolor=blue,
linkcolor=blue]{hyperref}

% ---------------------------------------------------------
% COMMANDES PERSONNELLES
% ---------------------------------------------------------
\newcommand\nomauteur{bdminasmorgul@protonmail.com}
\newcommand\dateTe{28.04.2026}
\newcommand\brancheTe{TS}

% ---------------------------------------------------------
% STYLE DE L'EXERCICE
% ---------------------------------------------------------
\pagestyle{headandfoot}
\firstpageheadrule
\runningheadrule

\firstpageheader{\brancheTe\ (\nomauteur)}{Élève : }{(\dateTe) \thepage/\numpages}
\runningheader{\brancheTe\ (\nomauteur)}{Élève : }{(\dateTe) \thepage/\numpages}

% Compteur en bas de page
\newcommand{\totalpointtts}{%
	\begin{tikzpicture}[remember picture, overlay]
		\node[anchor=south west, inner sep=0pt, outer sep=0pt]
		at ([xshift=0.5cm,yshift=0.5cm]current page.south west)
		{\rotatebox{90}{Total des points : \numpoints}};
	\end{tikzpicture}
}

\firstpagefooter{\totalpointtts}{}{}
\runningfooter{\totalpointtts}{}{}

% Réglages exam
\printanswers
\addpoints
\pointsinmargin
\bracketedpoints

% ---------------------------------------------------------
% LISTINGS (code)
% ---------------------------------------------------------
\usepackage{xcolor}
\usepackage{listings}

\lstdefinestyle{octavestyle}{
	language=Matlab,
	basicstyle=\ttfamily\small,
	keywordstyle=\color{blue},
	commentstyle=\color{green!50!black},
	stringstyle=\color{red},
	stepnumber=1,
	frame=single,
	breaklines=true,
	breakatwhitespace=true,
	tabsize=4
}
\usepackage{enumitem}
% Réglage global : toutes les listes enumerate (niveau 1) en a), b), c), ...
\setlist[enumerate,1]{label=\alph*), ref=\alph*)}
\newenvironment{explication}{%
	\par\noindent\textbf{Explication. }\itshape
}{\par\normalfont}


%\renewcommand{\thepartno}{\arabic{partno}}
%\renewcommand{\partlabel}{\thepartno)}
\begin{document}
	\unframedsolutions
	\pointsinmargin
	\bracketedpoints
	
\section*{Préparation TS PQ CFC ELMO,IELE,PELE}
\begin{questions}
\question L'énergie absorbée par un appareil est de \qty{34}{[kWh]}.\\ Il a fonctionné de \qty{8}{[h]} à \qty{20}{[h]} sous une tension de \qty{230}{[V]}.
\begin{parts}
	\part[3] Calculer l'intensité du courant dans l'appareil~?
	\part[3] Quelle est la valeur du disjoncteur de protection~?
	\part[3] Calculer la résistance de cet appareil totalement ohmique~?
\end{parts}


\begin{solution}
	Calcul de la durée de fonctionnement ($t$) :
	\begin{align*}
		t &= t_{fin} - t_{debut} \\[6pt]
		t &= \qty{20}{[h]} - \qty{8}{[h]} = \qty{12}{[h]}
	\end{align*}
	
	Calcul de la puissance absorbée ($P$) :
	\begin{align*}
		P &= \dfrac{W}{t} \\[6pt]
		P &= \dfrac{\qty{34}{[kWh]}}{\qty{12}{[h]}} = \qty{2,8333}{[kW]} = \qty{2833,33}{[W]}
	\end{align*}
	
	Calcul de l'intensité du courant ($I$) :
	\begin{align*}
		I &= \dfrac{P}{U} \\[6pt]
		I &= \dfrac{\qty{2833,33}{[W]}}{\qty{230}{[V]}} = \qty{12,3188}{[A]}
	\end{align*}
	
	Choix du disjoncteur :
	Le calibre nominal du disjoncteur ($I_n$) doit être supérieur à l'intensité nominale de l'appareil pour éviter les déclenchements intempestifs. Pour un courant de \qty{12,32}{[A]}, on choisit généralement un calibre standard de :
	\begin{align*}
		I_n &= \qty{13}{[A]}
	\end{align*}
	
	Calcul de la résistance de l'appareil ($R$) :
	\begin{align*}
		R &= \dfrac{U}{I} \\[6pt]
		R &= \dfrac{\qty{230}{[V]}}{\qty{12,3188}{[A]}} = \qty{18,6706}{[\Omega]}
	\end{align*}
	\vspace{0.5em}
	\needspace{20\baselineskip}
	\begin{verbatim}
		% Télécharger OCTAVE depuis https://wiki.octave.org/Using_Octave et l'installer
		% Puis copier-coller le code dans OCTAVE
		
		% Données de l'énoncé
		W_kwh = 34;         % Énergie absorbée [kWh]
		t_debut = 8;        % Heure de début [h]
		t_fin = 20;         % Heure de fin [h]
		U = 230;            % Tension d'alimentation [V]
		
		% 1. Calcul de la durée de fonctionnement
		t = t_fin - t_debut;
		printf("Durée de fonctionnement : t = %d [h]\n", t);
		
		% 2. Calcul de la puissance absorbée
		% Conversion de kWh en W (W_wh / t)
		P = (W_kwh * 1000) / t;
		printf("Puissance absorbée : P = %.2f [W]\n", P);
		
		% 3. Calcul de l'intensité du courant
		I = P / U;
		printf("Intensité du courant : I = %.4f [A]\n", I);
		
		% 4. Calcul de la résistance
		R = U / I;
		printf("Résistance de l'appareil : R = %.4f [Ohm]\n", R);
	\end{verbatim}
\end{solution}
\end{questions}
\end{document}